\documentclass[11pt,a4paper]{article}

\usepackage[margin=2.5cm]{geometry}
\usepackage{amsmath,amssymb,amsthm}
\usepackage{algorithm}
\usepackage{algpseudocode}
\usepackage{booktabs}
\usepackage{graphicx}
\usepackage{xcolor}
\usepackage{hyperref}

\hypersetup{colorlinks=true,linkcolor=blue!50!black,citecolor=blue!50!black,
            urlcolor=blue!50!black}

\newtheorem{proposition}{Proposition}
\newtheorem{remark}{Remark}

\newcommand{\xb}{\mathbf{x}^{b}}
\newcommand{\xa}{\mathbf{x}^{a}}
\newcommand{\Bmat}{\mathbf{B}}
\newcommand{\Lmat}{\mathbf{L}}
\newcommand{\Dmat}{\mathbf{D}}
\newcommand{\Hmat}{\mathbf{H}}
\newcommand{\Rmat}{\mathbf{R}}

\title{\textbf{Adaptive localization by observation assignment\\
for the modified-Cholesky ensemble Kalman filter}}

\author{Working draft}
\date{\today}

\begin{document}
\maketitle

\begin{abstract}
\noindent
Localization in ensemble data assimilation is normally parameterized by a
radius of influence, tuned once against a known truth and then held fixed. We
replace the radius by a different object: each model component is updated by
exactly one observation, chosen from the few that its neighbourhood reaches,
and the radius is whatever distance that choice implies. The choice is made by
minimizing a variational cost evaluated at a scalar local analysis, computed
from ensemble anomalies and observations alone, with no reference trajectory
and nothing withheld. Because the cost decomposes over components, the optimal
assignment is obtained exactly, in closed form, at a cost of microseconds per
evaluation. The resulting radius field adapts to the local observation density
without any parameter controlling it, and the observation selected is very
often not the nearest one: on a quasi-geostrophic testbed the second, third and
fourth candidates are chosen about as often as the first, because what decides
the choice is the ensemble correlation and not the distance. The assignment
determines the predecessor sets of a modified-Cholesky precision matrix, which
is assembled once per cycle and used for a single global analysis.
\end{abstract}

\section{The parameter nobody can pin down}

The analysis step of an ensemble Kalman filter combines a forecast ensemble
with observations. With an ensemble far smaller than the state dimension, the
sample covariance is rank deficient and carries spurious correlations between
distant variables; localization removes them. Every localization scheme depends
on a radius of influence, and the radius is normally chosen by sweeping it
against a known truth in a synthetic experiment and then freezing it.

Attempts to estimate the radius instead run into a difficulty that is not a
matter of tuning. A radius per model component is a free parameter with nothing
in the data to pin it down. In experiments preceding this work, on a system of
40 components, an oracle with full access to the truth reduced the analysis
error by 59\% on the cycle it optimized, and the same radius field made the
error 22\% \emph{worse} on the following cycle, performing worse than a random
shuffle of its own values. What such a search finds is the cycle's noise.

Grouping components reduces the count but introduces a grouping that must
itself be justified, and the radius within a group remains a free number.

\section{Assignment instead of radius}

We propose a different parameterization. Let the model state have components
$i = 1, \dots, n$ and let there be $m$ observations at known locations. For each
component, grow a neighbourhood until it contains $\kappa$ observations; call
them the \emph{candidates} of $i$ and write $\mathcal{C}(i)$ for the set. The
component is then updated by exactly one member of $\mathcal{C}(i)$, or by none
at all.

The decision variable is the assignment
%
\begin{equation}
  a : \{1, \dots, n\} \;\longrightarrow\; \mathcal{C}(i) \cup \{\varnothing\},
\end{equation}
%
and the radius of component $i$ is not chosen: it is the distance to $a(i)$.

Three properties follow immediately.

\paragraph{The radius adapts to the observation density with no parameter.}
Where observations are dense the neighbourhood stops growing early; where they
are sparse it grows further. On the testbed below, a regular lattice of spacing
four gives a mean radius of 2.7 and a maximum of 4, while a random network of
the same overall density gives a mean of 1.7 and a maximum of 5 --- the second
number is larger precisely because a random network leaves empty regions. No
setting controls this.

\paragraph{The nearest observation is not automatically the right one.}
What determines how much an observation tells us about a component is the
ensemble correlation between them, not the distance. In a flow with structure,
two points aligned along a current can be far better correlated than two
adjacent points across a front. This is why the neighbourhood keeps growing
past the first observation: the choice is made on the cost, and the cost sees
the correlation.

\paragraph{Abstaining is a legal choice.}
If none of the candidates correlates with the component, updating it from any
of them injects noise, which is what localization exists to prevent. Leaving
the component at its background value is one of the options, and the cost
selects it on its own.

\section{The cost}

For component $i$ assigned to observation $j$, the local analysis is the scalar
Kalman update
%
\begin{equation}
  x^{a}_{i} = x^{b}_{i} + \frac{\sigma_{ij}}{\sigma_{jj} + r_{j}}
              \big( y_{j} - x^{b}_{j} \big),
  \label{eq:local}
\end{equation}
%
with $\sigma_{ij}$ the ensemble covariance between component $i$ and the
observed site of $j$, $\sigma_{jj}$ the ensemble variance there, and $r_{j}$
the observation error variance. The cost is the variational one evaluated at
that analysis,
%
\begin{equation}
  J_{i}(a) \;=\;
  \underbrace{\frac{(x^{a}_{i} - x^{b}_{i})^{2}}{\sigma_{ii}}}_{\text{departure from the background}}
  \;+\;
  \underbrace{\frac{(y_{j} - x^{a}_{j})^{2}}{r_{j}}}_{\text{residual}} ,
  \qquad j = a(i),
  \label{eq:cost}
\end{equation}
%
and $J_{i} = (y_{j} - x^{b}_{j})^{2} / r_{j}$ when $i$ abstains, since the
background is kept and the departure term vanishes. The global score is
%
\begin{equation}
  J(a) \;=\; \frac{1}{n}\sum_{i=1}^{n} J_{i}(a).
  \label{eq:global}
\end{equation}

No reference trajectory enters \eqref{eq:cost}, and no observation is withheld.
Everything is a dot product over the $N$ ensemble members. An assignment that
moves the state a long way without justification pays in the first term; one
that leaves a residual it could have removed pays in the second.

\begin{proposition}[Exact minimization]
$J$ decomposes over components, so
$\arg\min_{a} J(a)$ is obtained by minimizing each $J_{i}$ independently over
$\mathcal{C}(i) \cup \{\varnothing\}$.
\end{proposition}

\begin{remark}
This is worth stating plainly because it removes a whole apparatus. There is no
combinatorial search here: the optimal assignment is a per-component
$\arg\min$ over $\kappa + 1$ precomputed numbers. Metaheuristics applied to
\eqref{eq:global} return exactly the same assignment --- in the experiments
below, tabu search at $2\times10^{5}$ evaluations and simulated annealing at
$1.6\times10^{4}$ both reproduce the greedy solution to every digit. The
decomposition is a feature of this cost, not of the implementation, and it is
lost as soon as a term coupling neighbouring components is introduced.
\end{remark}

\section{From assignment to analysis}

The assignment determines a radius per component, and therefore the predecessor
sets of the modified-Cholesky estimator \cite{bickel2008,nino2018}. For each
component $i$ with predecessors $P(i)$ given by its radius,
%
\begin{equation}
  \boldsymbol{\beta}_{i}
  = \big(\mathbf{X}_{P}^{\top}\mathbf{X}_{P} + \alpha\mathbf{I}\big)^{-1}
    \mathbf{X}_{P}^{\top}\,\Delta\mathbf{X}_{i},
  \qquad
  \Lmat_{ii} = 1, \quad \Lmat_{i,P(i)} = -\boldsymbol{\beta}_{i}^{\top},
\end{equation}
%
with $\Dmat_{ii}$ the inverse residual variance, and
$\widehat{\Bmat}^{-1} = \Lmat^{\top}\Dmat\Lmat$. The analysis solves
%
\begin{equation}
  \big(\widehat{\Bmat}^{-1} + \Hmat^{\top}\Rmat^{-1}\Hmat\big)\,
  \delta\mathbf{x} = \Hmat^{\top}\Rmat^{-1}\big(\mathbf{y} - \Hmat\xb\big).
\end{equation}
%
The precision is positive definite for any assignment, because it retains the
form $\Lmat^{\top}\Dmat\Lmat$; truncating an already-formed precision does not
preserve this, and in a companion measurement produced an indefinite matrix in
44 of 44 cases tested.

\begin{algorithm}[t]
\caption{One assimilation cycle}
\begin{algorithmic}[1]
\State propagate the ensemble; form anomalies $\Delta\mathbf{X}$
\State \textbf{for each component} grow its neighbourhood until it holds
       $\kappa$ observations
\State precompute $\sigma_{ii}$, $\sigma_{jj}$ and $\sigma_{ij}$ for every
       candidate pair
\State \textbf{for each component} take
       $a(i) = \arg\min J_{i}$ over $\mathcal{C}(i)\cup\{\varnothing\}$
\State read off the radius of every component from $a$
\State assemble $\widehat{\Bmat}^{-1}$ once and solve for the analysis
\end{algorithmic}
\end{algorithm}

\section{Numerical experiments}

The testbed is the 1.5-layer quasi-geostrophic model of Sakov and
Oke~\cite{sakov2008},
%
\begin{equation}
  q_t = -\psi_x - \varepsilon J(\psi,q) - A\triangle^{3}\psi + 2\pi\sin(2\pi y),
  \qquad q = \triangle\psi - F\psi,
\end{equation}
%
on the unit square with $F = 1600$, $\varepsilon = 10^{-5}$ and Dirichlet
boundaries, at $49 \times 49$ per field. The integration is carried on $q$ and
$\psi$ is recovered from it at every step, so $\psi$ is diagnostic and only $q$
is estimated: $n = 2401$. Ensemble members and the truth are drawn from
snapshots of a long run, following \cite{sakov2008}. The dissipation is ten
times the model default, which is what gives a stationary regime: at the
default the norm of $q$ grows by a factor of fourteen between $t=500$ and
$t=8000$ and there is no climatology to sample.

\subsection{What the assignment does}

With 144 observations on a lattice of spacing four and $\kappa = 4$:

\begin{table}[h]
\centering
\begin{tabular}{@{}lr@{}}
\toprule
Assignment & $J$ \\
\midrule
Every component abstains & 123.82 \\
Every component takes its nearest observation & 0.587 \\
Optimal assignment & \textbf{0.052} \\
Random assignment & 30.28 \\
\bottomrule
\end{tabular}
\caption{The cost separates good assignments from bad ones by orders of
magnitude, and the optimal assignment improves on taking the nearest
observation by a factor of eleven.}
\end{table}

The distribution of choices is the result worth emphasising. Of the 2401
components, 516 take their first candidate, 603 their second, 647 their third
and 635 their fourth. The candidates are ordered by distance, so the nearest
observation is chosen barely a fifth of the time: the ensemble correlation, not
the distance, is what decides.

\subsection{Cost}

Evaluating \eqref{eq:global} takes $1.6\times10^{-2}$ ms. Building the
candidate sets and precomputing the covariances is a fraction of a second per
cycle and is done once. For comparison, one evaluation of an objective that
requires forming the precision and solving the analysis system costs 0.18~s at
the same resolution, four orders of magnitude more; a radius large enough to
destroy the sparsity of the precision --- 51\% non-zero at radius 12 against
0.5\% at radius 1 --- pushes that to 1.07~s.

\section{Discussion}

The contribution is a reparameterization. The radius stops being a free number
and becomes a consequence of a choice among observations that are actually
present, which is what makes it estimable: a choice among $\kappa+1$
alternatives is constrained in a way that a positive real number is not.

The honest limitation is stated in the remark above. This cost decomposes, so
the optimization is trivial and no search is required. That is a virtue
operationally and a limitation scientifically: it means the assignment
formulation, as posed here, does not exercise the combinatorial structure it
appears to have. Coupling neighbouring components --- penalizing roughness in
the analysis increment, or capping how many components a single observation may
update --- breaks the decomposition and makes the problem genuinely
combinatorial. Whether that coupling improves the analysis, rather than merely
making the optimization harder, is a separate question and is left to
subsequent work.

\begin{thebibliography}{9}
\bibitem{bickel2008} P.~J. Bickel and E.~Levina.
  Regularized estimation of large covariance matrices.
  \emph{The Annals of Statistics}, 36(1):199--227, 2008.
\bibitem{nino2018} E.~D. Nino-Ruiz, A.~Sandu, and X.~Deng.
  An ensemble Kalman filter implementation based on modified Cholesky
  decomposition for inverse covariance matrix estimation.
  \emph{SIAM Journal on Scientific Computing}, 40(2):A867--A886, 2018.
\bibitem{nino2019} E.~D. Nino-Ruiz and L.~E. Morales-Retat.
  A Tabu Search implementation for adaptive localization in ensemble-based
  methods. \emph{Soft Computing}, 23:5519--5535, 2019.
\bibitem{sakov2008} P.~Sakov and P.~R. Oke.
  A deterministic formulation of the ensemble Kalman filter: an alternative to
  ensemble square root filters. \emph{Tellus A}, 60(2):361--371, 2008.
\end{thebibliography}

\end{document}
